\documentclass[10pt,a4paper]{amsart}
\usepackage[margin=19mm]{geometry}
\usepackage{amsmath,amssymb,mathtools}
\usepackage[scr=rsfs,cal=euler]{mathalfa}
\usepackage{xfrac}
\usepackage[unicode,hidelinks,bookmarks=false]{hyperref}
\newcommand{\ie}{i.e.\ }
\newcommand{\cf}{cf.\ }
\newcommand{\resp}{respectively\ }
\newcommand{\Subsection}{\subsection}
\providecommand{\nobreakdash}{\nobreak\hbox{-}}
\setlength{\emergencystretch}{2em}
\multlinegap0pt
\usepackage{amssymb}
\usepackage[shortlabels]{enumitem}
\newtheorem{theorem}{Theorem}
\newtheorem{proposition}{Proposition}
\newcommand{\N}{\mathbb{N}}
\newcommand{\calC}{\mathcal{C}}
\newcommand{\calS}{\mathcal{S}}
\newcommand{\calZ}{\mathcal{Z}}
\newcommand{\stau}{\overline\tau}
\title{Pluripotential theory on algebraic curves}
\author{Norm Levenberg, Sione Ma'u}
\date{Source: arXiv:2601.07639v2 (CC BY 4.0)}
\begin{document}
\maketitle

\noindent\emph{Source and licence.} Pluripotential theory on algebraic curves. Version arXiv:2601.07639v2 (CC BY 4.0), distributed by arXiv under the Creative Commons Attribution 4.0 International licence, http://creativecommons.org/licenses/by/4.0/, as recorded in the arXiv licence metadata for this exact version and shown on the abstract page https://arxiv.org/abs/2601.07639v2. Full licence notice: \texttt{licenses/arxiv-2601.07639-CC-BY-4.0.txt}.\par\smallskip

\noindent\textit{Editorial scope.} Counted unit: the article's proposition prop:13 (source0.tex lines 822--827). The article's complete original proof of it is delivered with it (source0.tex lines 829--889, one \texttt{proof} environment of 61 lines). No mathematics is written, repaired or re-derived: the statement and the proof are the article's own lines, in the article's own notation, and no retained line is substituted or re-flowed.  Retained uncounted as the source-local context the proof needs: lines 751--757; lines 781--786; lines 804--804.  Nothing was omitted from any retained span, so the package carries no omission record.  No retained line contains a citation, so the wrapper carries no bibliography and no bibliography entry.  No retained line contains a figure environment, an image-inclusion command or drawing code, so no image is delivered and the package declares no asset dependency.  Editorial formatting only, in addition: an AMS \texttt{amsart} preamble that restates the article's own declarations and macros used by the retained lines, namely the environments \texttt{theorem}, \texttt{proposition} on separate counters, together with the standard packages those lines need.  The retained environments print in this document as Theorem 1, Proposition 1 in order of appearance, whereas the article prints the same items with the numbering of its own full document.  The complete native source of the article as distributed in the arXiv e-print for this exact version is bundled unchanged as \texttt{sources/192520/source0.tex}.
\begin{theorem}\label{thm:ds=dc}
Let $K\subset A$ be compact, and suppose the directions $\lambda_1,\ldots,\lambda_d$ are labelled so that $T(K,\lambda_k)\geq T(K,\lambda_{k+1})$ for all $k\in\{1,\ldots, d-2\}$.  Then 
\begin{equation}\label{eqn:ds=dc}
\left(\prod_{k=1}^d  T(K,\calZ(k-1))\right)^{1/d} =  d_{\calS}(K) = d_{\calC}(K) = \left(\prod_{k=1}^d T(K,\lambda_k)\right)^{1/d},
\end{equation}
and hence $T(K,\calZ(k-1))=T(K,\lambda_k)$ for all $k\in\{1,\ldots,d\}$.
\end{theorem}

\begin{equation}\label{eqn:tp} \begin{aligned}
V_n=V_{N-1}\frac{V_{N}}{V_{N-1}}\cdots\frac{V_n}{V_{n-1}}&\leq V_{N-1} \frac{n!}{N!} \prod_{\nu=0}^{n-N}  \tau_{\nu+N}^{\alpha_{\nu+N}} \\ 
&   \\
&\leq V_{N-1} \frac{n!}{N!} \prod_{\nu=0}^{dm} (1+\epsilon)^{\alpha_{\nu+N}}\stau(\nu+N)^{\alpha_{\nu+N}} 
 \prod_{\nu=dm+1}^{dm+k} C^{\alpha_{\nu+N}} \\ 
\end{aligned}\end{equation}

\begin{equation}\label{eqn:ln} l_n=\frac{dm^2}{2}+O(m).\end{equation}

\begin{proposition} \label{prop:13}
Let $\Sigma_m:=\sum_{\nu=1}^m \nu$.  Then
\begin{equation}\label{eqn:TZ1}
\overline T(K,\calZ(1)) = \lim_{m\to\infty} \left(\prod_{\nu=1}^m \tau_{N+d\nu}^{\nu}\right)^{1/\Sigma_m} .
\end{equation}
\end{proposition}

\begin{proof}
We estimate as in \eqref{eqn:tp}, with $n=N-2 +dm$, $m\in\N$. (We will then let $m\to\infty$.)  We have 
\begin{equation}\label{eqn:prop415a}
V_n \leq V_{N-2}\frac{n!}{(N-2)!}\prod_{j=N-1}^n \tau_j^{\alpha_j}  = V_{N-2}\frac{n!}{(N-2)!} \prod_{k=0}^{d-1}\prod_{\nu=0}^m \tau_{N-1+d\nu+k}^{d-1+\nu}.
\end{equation}
 As $m\to\infty$ we have  $n\to\infty$, and therefore 
\begin{equation}\label{eqn:T1a}
\begin{aligned}
d(K)  \leq   \limsup_{m\to\infty}\left(\prod_{k=0}^{d-1}\prod_{\nu=0}^m \tau_{N-1+d\nu+k}^{d-1+\nu} \right)^{1/l_n}. 
\end{aligned}\end{equation}
(As before, the other factors on the right-hand side of \eqref{eqn:prop415a} go to $1$ in the limit.)  We will use \eqref{eqn:T1a} to show that
\begin{equation}\label{eqn:T1b}
\lim_{m\to\infty}\left(\prod_{\nu=0}^m \tau_{N+d\nu}^{d-1+\nu}\right)^{1/l_n} = \overline T(K,\calZ(1))^{1/d}.
\end{equation}
Let $\epsilon>0$.  Then for each $k$, $\tau_{N-1+d\nu+k}<(1+\epsilon)\overline T(K,\calZ(k))$ holds for all but finitely many $\nu$; we apply this to all terms where $k\neq 1$ on the right-hand side of \eqref{eqn:T1a}: 
$$\begin{aligned}
\limsup_{m\to\infty}\left(\prod_{k=0}^{d-1}\prod_{\nu=0}^m \tau_{N-1+d\nu+k}^{d-1+\nu} \right)^{1/l_n}
  & \leq \limsup_{m\to\infty}\left(\prod_{k\neq 1}\prod_{\nu=0}^m (1+\epsilon)\overline T(K,\calZ(k))^{d-1+\nu} \right)^{1/l_n} \\
& \qquad \times \limsup_{m\to\infty}\left(\prod_{\nu=0}^m \tau_{N+d\nu}^{d-1+\nu} \right)^{1/l_n} \\
& = \left((1+\epsilon)\prod_{k\neq 1} \overline T(K,\calZ(k)))\right)^{1/d} \\
&\qquad \times \limsup_{m\to\infty}\left(\prod_{\nu=0}^m \tau_{N+d\nu}^{d-1+\nu} \right)^{1/l_n}. \\
\end{aligned}$$ 
Then using Theorem \ref{thm:ds=dc}, 
$$\begin{aligned}
\left(\prod_{k=0}^{d-1} \overline T(K,\calZ(k)))\right)^{1/d}  =d(K)
& \leq \left((1+\epsilon)\prod_{k\neq 1} \overline T(K,\calZ(k)))\right)^{1/d} \\ 
&\qquad \times \limsup_{m\to\infty}\left(\prod_{\nu=0}^m \tau_{N+d\nu}^{d-1+\nu} \right)^{1/l_n}. 
\end{aligned}$$  Cancelling the common factors $\overline T(K,\calZ(k))$ for $k\neq 1$ on each side, and letting $\epsilon\to 0$, 
$$
\overline T(K,\calZ(1)))^{1/d} \leq \limsup_{m\to\infty}\left(\prod_{\nu=0}^m \tau_{N+d\nu}^{d-1+\nu} \right)^{1/l_n}.
$$
To get  a reverse inequality, let $\epsilon>0$.  Then $\tau_{N+dm}<(1+\epsilon)\overline T(K,\calZ(1))$ for all but finitely many $m\in\N$, so  
$$\begin{aligned}
\limsup_{m\to\infty}\left(\prod_{\nu=0}^m \tau_{N+d\nu}^{d-1+\nu} \right)^{1/l_n} &\leq \limsup_{m\to\infty}  \left(\prod_{\nu=0}^m ((1+\epsilon)\overline T(K,\calZ(1)))^{d-1+\nu}  \right)^{1/l_n} \\
& =    \limsup_{m\to\infty}\left( (1+\epsilon)\overline T(K,\calZ(1)) \right)^{\left(\sum_{\nu=0}^{m}(d-1-\nu)\right)/l_n }.
\end{aligned}$$ 
Using \eqref{eqn:ln}, we have  $\frac{1}{l_n}\sum_{\nu=0}^{m}(d-1-\nu)= \frac{m^2+O(m)}{dm^2+O(m)}\to \frac{1}{ d}$  as $m\to\infty$,  
 and we obtain  $$\limsup\limits_{m\to\infty}\left(\prod_{\nu=0}^m \tau_{N+d\nu}^{d-1+\nu} \right)^{1/l_n}\leq \left((1+\epsilon)\overline T(K,\calZ(k))\right)^{1/d}.$$  Letting $\epsilon\to 0$ we get the reverse inequality; altogether, 
$$
\overline T(K,\calZ(1)))^{1/d} = \limsup_{m\to\infty}\left(\prod_{\nu=0}^m \tau_{N+d\nu}^{d-1+\nu} \right)^{1/l_n}.
$$
This is almost \eqref{eqn:T1b}, except that we need limsup replaced by lim.  

For the sake of obtaining a contradiction, suppose the limsup is not a limit. Then there exists $\delta>0$ and 
a subsequence $\{m_j\}$ in the parameter $m$ such that 
$$
\overline T(K,\calZ(1)))^{1/d}-\delta = \lim_{k\to\infty} \left(\prod_{\nu=0}^{m_j} \tau_{N+d\nu}^{d-1+\nu} \right)^{1/l_{n_j}}.
$$
But then, by similar reasoning as above with the limit formula $d(K)=\lim_{j\to\infty} (V_{n_j})^{1/l_{n_j}}$ (where $n_j=N-2+dm_j$),  we obtain for any $\epsilon>0$, 
$$\begin{aligned}
\left(\prod_{k=0}^{d-1} \overline T(K,\calZ(k))\right)^{1/d}  
 \leq \left((1+\epsilon)\prod_{k\neq 1} \overline T(K,\calZ(k)))\right)^{1/d}  
\left(\overline T(K,\calZ(1)))^{1/d}-\delta\right)
\end{aligned}$$
and therefore $\overline T(K,\calZ(1))^{\frac{1}{d}}\leq (1+\epsilon)^{1-\frac{1}{d}}(\overline T(K,\calZ(1)))^{\frac{1}{d}}-\delta )$.  
Letting $\epsilon\to 0$ gives a contradiction.  Thus \eqref{eqn:T1b} holds.  

To finish the proof, we must modify \eqref{eqn:T1b} to obtain \eqref{eqn:TZ1}.  First, we may replace  $\tau_{N+d\nu}^{d-1+\nu}$ in the product by  $\tau_{N+d\nu}^{\nu}$. This amounts to dividing out $\tau_{N+d\nu}^{d-1}$ for each $\nu=1,\ldots,m$.  If we estimate all of these $\tau_{N+d\nu}$ quantities by a uniform constant $C$ then  \eqref{eqn:T1b} is off by an estimated $C^{m(d-1)}$  relative to \eqref{eqn:TZ1}; the $l_n$-th root of this goes to 1 as $m\to\infty$  (recall $l_n=dm^2/2+O(m)$).  

Finally,  $\Sigma_m= \frac{m^2}{2}+O(m)$, so $\frac{l_n}{\Sigma_m}\to d$.  Replacing the $l_n$-th root in the formula by the $\Sigma_m$-th root,  the new limit as $m\to\infty$ is the $d$-th power of the original limit.  Hence \eqref{eqn:TZ1} holds.  
\end{proof}

\end{document}
